\section{Benchmark Methodology}
\label{app:benchmark}

This appendix describes the experimental procedures used for the performance evaluation.

\subsection*{Timing Methodology}

All time measurements use Tcl's built-in \texttt{time} command, which executes a script
for a specified number of iterations and returns the average execution time per iteration
in microseconds. To ensure fair comparisons, a \texttt{profile} wrapper discards the first
execution (allowing Tcl's bytecode compiler to optimize the procedure) before measuring:

\begin{lstlisting}[language=Tcl, caption={Timing wrapper discarding the first run}]
proc profile {body times} {
    uplevel 1 $body; # discard first run to compile tcl procs
    uplevel 1 [list time $body $times]
}
\end{lstlisting}

Each benchmark invocation runs with 1,000 iterations for object creation, getter, and
setter tests. Class declaration benchmarks also use 1,000 iterations but measure the time
to declare complete classes with multiple fields and methods.

\subsection*{Test Classes}

The benchmark suite uses a \texttt{Point} class with five fields of mixed types to
represent realistic object complexity:

\begin{lstlisting}[language=Tcl, caption={VOO \texttt{Point} implementation}]
voo::class VooPoint {
    public {
        double_t x 0.0
        double_t y 0.0
        string_t name "point"
        int_t id 0
        bool_t active 1
    }

    method distance {} {
        set dx [get.x $this]
        set dy [get.y $this]
        return [expr {sqrt($dx * $dx + $dy * $dy)}]
    }
}
\end{lstlisting}

\begin{lstlisting}[language=Tcl, caption={TclOO \texttt{Point} implementation}]
oo::class create TclooPoint {
    variable x 0.0
    variable y 0.0
    variable name "point"
    variable id 0
    variable active 1

    constructor {x_ y_ name_ id_ active_} {
        set x $x_; set y $y_; set name $name_; set id $id_; set active $active_
    }

    method getX {} { return $x }
    method setX {value} { set x $value }
    method distance {} { return [expr {sqrt($x * $x + $y * $y)}] }
}
\end{lstlisting}

\begin{lstlisting}[language=Tcl, caption={Itcl \texttt{Point} implementation}]
itcl::class ItclPoint {
    public variable x 0.0
    public variable y 0.0
    public variable name "point"
    public variable id 0
    public variable active 1

    constructor {x_ y_ name_ id_ active_} {
        set x $x_; set y $y_; set name $name_; set id $id_; set active $active_
    }

    method getX {} { return $x }
    method setX {value} { set x $value }
    method distance {} { return [expr {sqrt($x * $x + $y * $y)}] }
}
\end{lstlisting}

The VOO~C++ implementation is loaded from a pre-compiled shared library:

\begin{lstlisting}[language=Tcl, caption={Loading and using the VOO C++ Point class}]
# VOO C++ --- load compiled shared library built from point.cpp
load <path/to/libvoopoint_cpp.so> Point

set obj [CppVooPoint::new 1.0 2.0 "test" 1 1]
set obj [CppVooPoint::new()]                    ;# default constructor
set x   [CppVooPoint::get.x $obj]              ;# getter --- object by value
CppVooPoint::set.x obj 3.14                    ;# setter --- variable by name
\end{lstlisting}

\noindent\textbf{Note on Class Declaration Performance:} VOO~C++ class declarations
happen entirely at C++ compile time; there is no runtime class registration cost to
measure. Package loading time (the one-time cost of the \texttt{load} command) is excluded
from all benchmark comparisons because it is an initialization step analogous to
\texttt{source}-ing the VOO Tcl framework or loading Itcl.

\subsection*{Memory Measurement}

Memory measurements were performed separately using a dedicated memory benchmark script.
Each scenario runs in an isolated \texttt{tclsh} process. The process resident set size
(RES) is recorded via \texttt{htop} after creating 100,000 objects of each class type.

\subsection*{Benchmark Categories}

The benchmark suite measures five distinct performance metrics:
\begin{enumerate}
  \item \textbf{Object Creation (Explicit Values):} Instantiate objects with all field
        values explicitly provided to the constructor.
  \item \textbf{Object Creation (Default Values):} Instantiate objects using the
        no-argument constructor with all default field values.
  \item \textbf{Setter Performance:} Modify a single field value through the setter.
  \item \textbf{Getter Performance:} Retrieve a single field value through the getter.
  \item \textbf{Class Declaration Performance:} Declare a complete class with fields and
        methods (one-time initialization cost).
\end{enumerate}
